% =============================================================================
%  Paper_V3.tex
%  Hybrid Neural Force Correction for the Unrestricted Three-Body Problem
%  Aarush Gupta  |  ST Yau High School Science Award (Asia)  |  2026
%
%  OVERLEAF INSTRUCTIONS:
%  1. Paste this file as your main .tex file
%  2. Paste V3_bib.bib content into a new file called "V3_bib.bib" in Overleaf
%  3. Upload Declaration_of_academic_integrity.pdf to Overleaf
%  4. Upload all 25 figure files listed at the top of the Results section
% =============================================================================

\documentclass[12pt]{article}

% ── Page geometry (ST Yau: A4, 2.5cm all sides) ───────────────────────────────
\usepackage[a4paper,
            left=2.5cm, right=2.5cm,
            top=2.5cm,  bottom=2.5cm]{geometry}

% ── Core packages ─────────────────────────────────────────────────────────────
\usepackage{graphicx}
\usepackage{amsmath}
\usepackage{amssymb}
\usepackage{float}
\usepackage{natbib}
\usepackage{setspace}       % line spacing
\usepackage{booktabs}       % professional tables
\usepackage{array}          % table column formatting
\usepackage{caption}        % caption formatting
\usepackage{pdfpages}       % include Academic Honesty PDF
\usepackage{fancyhdr}       % headers and footers
\usepackage{titlesec}       % section formatting
\usepackage{parskip}        % paragraph spacing
\usepackage{tabularx}

% ── Clickable TOC and cross-references ────────────────────────────────────────
\usepackage[
  colorlinks = true,
  linkcolor  = blue,
  citecolor  = blue,
  urlcolor   = blue,
  bookmarks  = true,
]{hyperref}

% ── Line spacing (ST Yau: 1.5) ────────────────────────────────────────────────
\onehalfspacing

% ── Header / footer ───────────────────────────────────────────────────────────
\pagestyle{fancy}
\fancyhf{}
\rhead{\small Hybrid Neural Force Correction for the Unrestricted Three-Body Problem}
\lhead{\small Aarush Gupta}
\cfoot{\thepage}
\renewcommand{\headrulewidth}{0.4pt}

% ── Caption formatting ────────────────────────────────────────────────────────
\captionsetup{font=small, labelfont=bf, width=0.95\linewidth}

% =============================================================================
\begin{document}
% =============================================================================

% ── COVER PAGE ────────────────────────────────────────────────────────────────
\begin{titlepage}
\thispagestyle{empty}

\begin{flushleft}
\textbf{Research Report} \hfill \textbf{2026 S.T. Yau High School Science Award (Asia)}
\end{flushleft}

\vspace{1.2cm}

\begin{center}
{\Large \textbf{2026 S.T. Yau High School Science Award (Asia)}}\\[0.3cm]
{\Large \textbf{Research Report}}\\[1.2cm]
\end{center}

\noindent\textbf{The Team}\\[0.4cm]

\noindent\textbf{Registration Number:} [Registration Number]\\[0.5cm]

\noindent\textbf{Name of team member:} Aarush Gupta\\
\textbf{School:} UWCSEA East Campus\\
\textbf{Country:} Singapore\\[0.8cm]

\noindent\textbf{Name of supervising teacher:} [Supervisor Name]\\
\textbf{Job Title:} [Supervisor Job Title]\\
\textbf{School/Institution:} [Supervisor School/Institution]\\
\textbf{Country:} [Supervisor Country]\\[1.2cm]

\noindent\textbf{Title of Research Report}\\[0.3cm]
\begin{center}
{\Large \textbf{Hybrid Neural Force Correction for the\\[0.2cm]
Unrestricted Three-Body Problem}}
\end{center}

\vfill

\begin{center}
\textbf{Date}\\
17 August 2026
\end{center}

\end{titlepage}

% ── TITLE, AUTHORS, ABSTRACT, KEYWORDS (Roman numeral page i) ─────────────────
\newpage
\pagenumbering{roman}
\setcounter{page}{1}
\thispagestyle{plain}

\begin{center}
  {\Large \textbf{Hybrid Neural Force Correction for the\\[0.2cm]
   Unrestricted Three-Body Problem}}\\[0.8cm]
  {\large Aarush Gupta}\\[0.3cm]
  {\small [UWCSEA], [Singapore]}\\[1cm]
\end{center}

\phantomsection
\section*{Abstract}
\addcontentsline{toc}{section}{Abstract}

We present SIMON, a hybrid integrator combining a second-order leapfrog scheme with a learned scalar pairwise force correction and two physical constraints: an analytic force direction that preserves Newtonian radial symmetry, and adaptive sub-stepping that resolves close gravitational encounters. Through a controlled four-run ablation study---removing each component in isolation---we show that long-horizon orbital stability over 100 years depends on these structural constraints: adaptive sub-stepping eliminates the timestep-frontier ejection failure (105$\times$ worst-case frontier error
reduction), while the learned-direction baseline demonstrates, for the tested initial condition, that replacing analytic direction with a neural approximation produces unbounded trajectories. The neural network correction itself provides an $8.5\%$ accuracy improvement without contributing to stability. Evaluated across seven distinct three-body initial conditions spanning near-circular to tight-binary dynamical regimes, SIMON maintains bounded dynamics for all stable configurations tested (mean divergence rate $\lambda = 0.092 \pm 0.068\ \mathrm{yr}^{-1}$), and agrees with the reference integrator on the occurrence of ejection in the three
dynamically unstable configurations; post-ejection separations differ, as expected, because
chaotic divergence amplifies integrator differences after the system becomes unbound. Energy drift scales as $dt^2$, consistent with second-order leapfrog theory, with negligible contribution from the neural correction. At the operational timestep ($dt = 0.04\ \mathrm{yr}$), SIMON achieves a $1.12\times$ speedup over \texttt{ias15} with simulation performance invariant across five independent training seeds. These results demonstrate that long-horizon stability in hybrid neural integrators for chaotic gravitational systems requires structurally enforced physical constraints rather than learned approximations alone.

\vspace{0.5cm}
\noindent\textbf{Keywords:} Three-Body Problem; Neural Networks; Hybrid Integrator;
Leapfrog Integration; Adaptive Sub-Stepping; Chaotic Dynamics; Gravitational N-Body Simulation; Lyapunov Divergence

% ── ACKNOWLEDGEMENT (page ii) ─────────────────────────────────────────────────
\newpage
\thispagestyle{plain}

\phantomsection
\section*{Acknowledgement}
\addcontentsline{toc}{section}{Acknowledgement}

% [Student fills in acknowledgement here]
\vspace{3cm}
\noindent [Acknowledgement text to be completed by author.]

% ── ACADEMIC HONESTY AND INTEGRITY (page iii) ─────────────────────────────────
%\newpage
% Include the signed Academic Honesty PDF.
% Upload "Declaration_of_academic_integrity.pdf" to your Overleaf project.
%\includepdf[pages={Declaration_of_academic_integrity.pdf}]
%\includepdf[pages={1}]{Declaration_of_academic_integrity.pdf}

% ── ACADEMIC HONESTY AND INTEGRITY (page iii) ─────────────────────────────────
\newpage
\phantomsection
\addcontentsline{toc}{section}{Commitments on Academic Honesty and Integrity}

% Include the one-page ST Yau Academic Integrity Statement.
% Upload the signed/scanned file to Overleaf with this exact filename:
% Declaration_of_academic_integrity.pdf
\includepdf[
  pages=1,
  pagecommand={\thispagestyle{plain}},
  fitpaper=true
]{Declaration_of_academic_integrity.pdf}

%===========================================================
% ── TABLE OF CONTENTS (pages iv-v) ────────────────────────────────────────────
\newpage
\renewcommand{\contentsname}{Table of Contents}
\tableofcontents

% =============================================================================
% MAIN REPORT — Arabic numerals
% =============================================================================
\newpage
\pagenumbering{arabic}
\setcounter{page}{1}

% =============================================================================
\section{Introduction}
\label{sec:intro}
% =============================================================================

The Three-Body Problem---the gravitational dynamics of three mutually interacting massive bodies---is a classical problem in physics because its trajectories are generally chaotic and highly sensitive to initial conditions. Although exact analytical solutions exist for two-body systems, long-horizon three-body prediction must usually be performed numerically. This remains important in astrophysical applications such as exoplanet stability, stellar cluster dynamics, and spacecraft trajectory design \cite{gardner2006james}.

High-accuracy numerical integrators such as \texttt{ias15} provide reliable reference solutions, but their step-by-step structure makes long simulations computationally expensive \cite{Rein_2014}. Neural or hybrid methods offer a possible route to faster simulation \cite{ulibarrena2024hybrid}, but unconstrained learned force models can lose physical stability during close encounters, where small position or force errors are strongly amplified \cite{Dehnen2001}. The central issue is therefore not simply whether a neural network can approximate a trajectory, but whether it can do so while preserving the physical structure needed for stable long-horizon dynamics.

This study addresses the following research question:
\begin{quote}
\emph{How closely can physically constrained neural network force models approximate the long-term stability of traditional numerical integrators in three-body simulations, while reducing computational cost without sacrificing qualitative dynamical stability?}
\end{quote}

To answer this question, we present \texttt{SIMON}, a hybrid integrator that combines a second-order leapfrog scheme with a learned scalar pairwise force correction, analytically enforced Newtonian force direction, and adaptive sub-stepping during close encounters. The method is evaluated through controlled ablations that isolate which components are responsible for stability, accuracy, physical consistency, and computational performance.

The key contributions of this work are:
\begin{enumerate}
    \item \textbf{A stability decomposition of hybrid neural integration:} controlled ablations show that adaptive sub-stepping and analytic force direction are the main stability mechanisms, while the neural correction primarily improves accuracy.

    \item \textbf{A generalisation test across dynamical regimes:} \texttt{SIMON} is evaluated on seven distinct initial conditions, including stable, tight-binary, hierarchical, and physically unstable configurations.

    \item \textbf{A physical-consistency and reproducibility analysis:} energy drift is shown to scale as $dt^2$, consistent with leapfrog theory, and results remain stable across five independent training seeds.

    \item \textbf{A speed--accuracy frontier:} at the operational timestep, \texttt{SIMON} achieves a modest computational-cost reduction while preserving bounded dynamics across all tested stable configurations.
\end{enumerate}

The paper is structured as follows.
Section~\ref{sec:lit_review} reviews the physics background of three-body dynamics,
chaos, and long-horizon prediction.
Section~\ref{sec:comp_modeling} discusses computational modelling, learned simulators,
and the research gap addressed by \texttt{SIMON}.
Section~\ref{sec:methodology} defines the \texttt{SIMON} architecture, training procedure,
evaluation strategy, and scope.
Section~\ref{sec:results} presents experimental results across four themes:
stability decomposition, generalisation, physical consistency, and computational
performance.
Section~\ref{sec:discussion} discusses implications, limitations, and future
directions.

% =============================================================================
\section{Literature Review}
\label{sec:lit_review}

The literature relevant to this project lies at the intersection of celestial mechanics,
numerical integration, and scientific machine learning. Classical mechanics explains why
the unrestricted three-body problem is sensitive to initial conditions and difficult to solve
analytically, while numerical-analysis literature shows that long-horizon behaviour depends
strongly on whether an algorithm preserves geometric structure and handles close encounters
robustly \cite{hairer2006geometric,Rein_2015,Rein_2014,boekholt2015reliability}. More recent
machine-learning work shows that neural networks can act as learned simulators or learned
Hamiltonian surrogates, but also that their success depends on the physical inductive biases
built into the architecture \cite{chen2018neuralode,battaglia2016interaction,greydanus2019hamiltonian,gruver2022deconstructing,horn2025generalized}. This section gives the physics background needed to motivate the computational gap addressed by \texttt{SIMON}.

\subsection{From Two-Body Solvability to Three-Body Non-Integrability}

The Two-Body Problem is analytically tractable because the mutual gravitational attraction
between two isolated point masses can be reduced to an equivalent one-body problem in a
central potential. Newton's law of universal gravitation gives the magnitude of the force as
\begin{equation}
F = G\frac{m_1m_2}{r^2},
\end{equation}
where $G$ is the gravitational constant, $m_1$ and $m_2$ are the two masses, and $r$ is their
separation. This framework explains the classical Keplerian orbits used to model systems
such as planets, moons, and satellites \cite{murray1999solar}.

The unrestricted Three-Body Problem is fundamentally different. Each body is accelerated
by the other two at the same time, so the equations of motion are coupled and nonlinear.
For bodies with positions $\mathbf r_i$ and masses $m_i$, the Newtonian equation of motion
can be written compactly as
\begin{equation}
\ddot{\mathbf r}_i =
G\sum_{j\ne i} m_j
\frac{\mathbf r_j-\mathbf r_i}{|\mathbf r_j-\mathbf r_i|^3},
\qquad i=1,2,3.
\end{equation}
Unlike the two-body case, this system has no general closed-form solution. Its behaviour
depends sensitively on mass ratios, separations, and velocities, making numerical simulation
the standard approach for studying generic three-body dynamics \cite{murray1999solar}.

\subsection{Chaos and Long-Horizon Prediction}
\label{sec:chaos}

A central difficulty of the Three-Body Problem is chaos. Small differences in initial position
or velocity can grow rapidly, causing two initially similar trajectories to diverge over long
times. One common way to describe this behaviour is the Lyapunov exponent $\lambda$, where
an initial separation $\delta(0)$ grows approximately as
\begin{equation}
\delta(t) \approx \delta(0)e^{\lambda t}.
\end{equation}
A positive value of $\lambda$ indicates exponential sensitivity to initial conditions \cite{ott2002chaos}.
This idea traces back to Poincar\'e's analysis of the three-body problem, which showed that
gravitational systems can exhibit intricate, non-repeating trajectories rather than simple
closed-form motion \cite{poincare1890sur}.

For this project, chaos has an important methodological consequence. In a chaotic system,
long-term pointwise trajectory agreement is not always the right standard for judging a
simulator, because even small numerical differences can produce phase drift. Instead, a
useful evaluation must also consider qualitative dynamical behaviour: whether a system
remains bounded, whether ejection occurs, whether energy drift is physically explainable,
and whether divergence grows at a controlled rate. This motivates the use of divergence-rate
proxies and boundedness checks later in this paper.

\subsection{Restricted and Unrestricted Three-Body Problems}

The Restricted Three-Body Problem is a simplified version in which two massive bodies
influence a third body whose mass is small enough to neglect. This formulation led to
important results such as Euler's and Lagrange's special solutions and the five Lagrange
points \cite{bistafa2021euler,lagrange1772essai}. These ideas are important in practical mission
design; for example, the James Webb Space Telescope operates near the Sun--Earth $L_2$
region \cite{gardner2006james}.

However, the present work focuses on the unrestricted problem, where all three bodies exert
mutual gravitational forces. This case is more directly relevant to testing neural-assisted
integrators because the system does not reduce to a fixed background potential and can
include strong pairwise encounters. The unrestricted setting therefore provides a stricter
test of whether a learned force model can remain dynamically usable when all interactions
are active.

The key implication for this project is that unrestricted three-body dynamics combine three
features that make long-horizon simulation difficult: nonlinear mutual forcing, sensitivity
to initial conditions, and close encounters where the Newtonian force changes rapidly with
separation. These features make purely pointwise long-term trajectory matching unreliable
and motivate the computational question addressed next: how can a hybrid neural integrator
reduce computational cost while preserving the physical structure needed for stable qualitative dynamics?
% =============================================================================
\section{Computational Modelling}
\label{sec:comp_modeling}

Since unrestricted three-body dynamics do not have a general closed-form solution,
their evolution must be computed numerically. A traditional integrator advances the
system step by step, so long simulations require many force evaluations, and close
encounters can amplify small numerical errors into large trajectory differences. High-accuracy
tools such as \texttt{REBOUND} and \texttt{ias15} are therefore essential benchmarks for
this project, because they show what reliable long-horizon gravitational integration
should look like \cite{Rein2012,Rein_2014}. At the same time, numerical-analysis
literature shows that long-time reliability is not determined only by local truncation
error: preserving geometric structure and resolving close interactions are central to
physical reliability \cite{hairer2006geometric,boekholt2015reliability}.

Neural networks offer a possible route to reducing this computational burden, but the
role of the network must be defined carefully. Some learned-surrogate approaches train
a model to map from an initial state directly to a future state, thereby skipping some
or all intermediate integration steps. \texttt{SIMON} does not use this direct
state-to-future-state strategy. It remains a timestep-based hybrid integrator, using
a neural network only as a local scalar force-correction mechanism during close
encounters. The aim is therefore not to replace numerical integration entirely, but to
test whether a learned component can reduce computational cost while preserving the
physical structure needed for stable long-horizon rollouts.

This distinction is important because close encounters are the most fragile part of the
problem. Under Newtonian gravity, the force between two bodies separated by distance
$r$ is proportional to $1/r^2$, so the force magnitude and its gradient change rapidly
when $r$ becomes small. A small learned error in position, force magnitude, or force
direction can then produce a large artificial impulse, causing non-physical energy
change or ejection-like behaviour \cite{Dehnen2001}. The central modelling question is
therefore not simply whether a neural network can approximate a three-body trajectory,
but which parts of the Newtonian force law must remain exact for a learned integrator
to remain dynamically usable.

\subsection{Learned Simulators and Research Gap}

Prior work gives three important lessons for this project. First, classical integrators
provide the reliability standard: \texttt{WHFast} demonstrates the value of fast
symplectic long-term orbit integration, while \texttt{ias15} provides a high-order
adaptive Gau\ss--Radau reference method for gravitational dynamics
\cite{Rein_2015,Rein_2014}. Work on the reliability of $N$-body simulations further
shows that chaotic three-body scattering can require very high numerical precision to
obtain converged individual trajectories \cite{boekholt2015reliability}. This makes
close encounters and long-horizon agreement especially demanding tests for any hybrid
method.

Second, learned simulators show that neural networks can approximate dynamical systems,
but they do not by themselves guarantee physical reliability. Neural ODEs learn a
continuous-time derivative field \cite{chen2018neuralode}, while Interaction Networks
and graph-network simulators learn object interactions and roll out multi-body systems
over time \cite{battaglia2016interaction,sanchez2020learning}. In celestial mechanics,
Breen \emph{et al.} showed that deep networks trained on high-precision three-body
solutions can approximate bounded-time three-body evolution at fixed cost
\cite{breen2020newton}. Physics-informed neural networks take another route by imposing
governing equations through the loss function \cite{cuomo2022scientific}, including
recent applications to the three-body problem \cite{pereira2025advancing}. These works
show the promise of learned dynamics, but they do not isolate which components of the
Newtonian force law should be learned and which should remain analytically enforced.

Third, structure-preserving machine learning shows that adding physical inductive bias
can improve extrapolation. Hamiltonian and Lagrangian Neural Networks encode mechanical
structure through energy-based formulations \cite{greydanus2019hamiltonian,cranmer2020lagrangian},
while symplectic or Hamiltonian-inspired architectures aim to preserve geometric
properties of the underlying dynamics \cite{zhong2020symoden,jin2020sympnets,tong2021symplectic,david2023symplectic,horn2025generalized}. However, later analysis also shows that
not every physics-inspired bias contributes equally; some apparent Hamiltonian-network
advantages may arise from other inductive biases such as second-order structure
\cite{gruver2022deconstructing}. This supports a more targeted question: in a chaotic
gravitational system, which physical constraints are actually necessary for stable
long-horizon rollouts?

This question is especially important because exact pointwise agreement over long
horizons is often unrealistic in chaotic systems. A useful learned model should remain
in the correct qualitative dynamical regime, rather than only minimising short-horizon
rollout error \cite{linot2022stabilized,margazoglou2023stability}. For the three-body
problem, this means preserving bounded motion when the reference solution remains
bounded, identifying ejection when the reference system becomes unstable, and keeping
energy drift physically interpretable.

The gap addressed by \texttt{SIMON} is therefore specific: prior work shows that learned
simulators can be fast and that physical inductive biases can improve generalisation, but
it does not clearly test whether a neural force model should learn the full gravitational
force vector or whether Newtonian radial direction should remain analytic while learning
is restricted to a lower-dimensional scalar correction. Nor does it usually separate the
stabilising effect of close-encounter timestep refinement from the effect of the neural
model itself. \texttt{SIMON} is designed around this gap. It combines a leapfrog backbone
with analytically enforced radial force direction, adaptive sub-stepping during close
encounters, and a neural network limited to a scalar pairwise force correction. Its
novelty is methodological: through controlled ablation, it tests which physical
constraints are necessary for dynamically usable long-horizon three-body rollouts.

% =============================================================================
\section{SIMON: Architecture and Methodology}
\label{sec:methodology}

\subsection{Method Overview and Reference Solver}
\label{research_obj}

Building on the research gap identified in Section~\ref{sec:comp_modeling}, this section
defines the hybrid architecture used to test whether a neural force correction can be made
dynamically usable in the unrestricted Three-Body Problem. \texttt{SIMON} is a timestep-based
integrator: it does not predict the final state directly from the initial state. Instead, it
evolves the system step by step using a leapfrog backbone, while using a neural network only
as a local scalar correction to the pairwise force magnitude during close encounters.

The reference solution throughout this paper is generated using \texttt{REBOUND}
\citep{Rein2012}, specifically the \texttt{ias15} integrator \citep{Rein_2014}. \texttt{ias15}
is a high-order adaptive Gau\ss--Radau method designed for gravitational dynamics and is
used here as the benchmark for long-horizon trajectory comparison, boundedness/ejection
classification, and runtime comparison. \texttt{SIMON} is not intended to replace
\texttt{ias15} as a universal reference solver. Its purpose is to test a narrower methodological
question: which physical constraints must be enforced for a neural-assisted integrator to
preserve qualitative long-horizon dynamics?

The method has three core design choices:
\begin{enumerate}
    \item a second-order leapfrog update provides the integration backbone;
    \item Newtonian radial force direction is computed analytically at every step;
    \item the neural network learns only a scalar correction to the softened pairwise force
    magnitude during close encounters.
\end{enumerate}
This separation between learned and analytic components allows the ablation study in
Section~5.1 to identify which parts of the method are responsible for stability, accuracy,
and computational performance.

\subsection{Scalar Force-Correction Model}

For bodies with positions $\mathbf{x}_i$, velocities $\mathbf{v}_i$, and masses $m_i$, pairwise
interactions are computed using the displacement
\[
\mathbf{r}_{ij}=\mathbf{x}_j-\mathbf{x}_i,\qquad
r_{ij}=\|\mathbf{r}_{ij}\|.
\]
The exact Newtonian pairwise force is
\[
\mathbf{F}_{ij}^{\mathrm{Newt}}
=
G\frac{m_i m_j}{r_{ij}^3}\mathbf{r}_{ij},
\]
where the vector direction is determined entirely by the geometry of the two bodies
\cite{murray1999solar}. This exact radial direction is preserved in \texttt{SIMON}; the neural
network is never allowed to learn or rotate the force direction.

Close encounters are numerically difficult because the Newtonian force magnitude changes
rapidly as $r_{ij}$ becomes small. To obtain a smooth local correction target, \texttt{SIMON}
first computes a softened pairwise force using
\[
\tilde r_{ij}=\sqrt{r_{ij}^2+\epsilon^2},
\qquad
\mathbf{F}_{ij}^{\mathrm{soft}}
=
G\frac{m_i m_j}{\tilde r_{ij}^{3}}\mathbf{r}_{ij}.
\]
The network then predicts a scalar correction factor $c_{ij}$ so that
\[
\mathbf{F}_{ij}^{\mathrm{SIMON}}
=
c_{ij}\mathbf{F}_{ij}^{\mathrm{soft}}.
\]
Equivalently, the ideal correction satisfies
\[
c_{ij}
=
\left(\frac{\tilde r_{ij}}{r_{ij}}\right)^3,
\qquad
\log c_{ij}
=
3\log\left(\frac{\tilde r_{ij}}{r_{ij}}\right).
\]
The model is therefore trained to predict $\log c_{ij}$ from the feature vector
\[
z_{ij}=[\log \tilde r_{ij}, \log m_i, \log m_j].
\]
This keeps the learned task one-dimensional in the output while preserving the exact
Newtonian force direction. The architecture is a small multilayer perceptron,
\[
3 \rightarrow 32 \rightarrow 32 \rightarrow 32 \rightarrow 1,
\]
with SiLU activations. This compact design is sufficient because the learned function is a
smooth scalar correction, not a full vector-valued force law.

\subsection{Hybrid Integration and Adaptive Sub-Stepping}
\label{sec:adaptive}

At each macro timestep, \texttt{SIMON} computes all pairwise forces and updates the state
using a velocity-Verlet/leapfrog step. The force calculation is divided into three regimes:

\begin{enumerate}
    \item \textbf{Far field:} if $r_{ij}>500\epsilon=0.15\,\mathrm{AU}$, the NN is not invoked and
    the pairwise force is computed using the exact Newtonian expression.
    \item \textbf{Close encounter:} if $r_{ij}<0.15\,\mathrm{AU}$, the NN predicts the scalar
    correction $c_{ij}$ applied to the softened force.
    \item \textbf{Safety fallback:} if the encounter is extremely close
    $(\tilde r_{ij}<5\times10^{-4}\,\mathrm{AU})$, or if the predicted correction lies outside
    the gate $0.2\leq c_{ij}\leq 5.0$, the NN output is rejected and the exact Newtonian
    force is used instead.
\end{enumerate}

This design follows the principle of hybrid scientific machine learning: learned components
are used only where they are intended to help, while physically unsafe outputs are replaced
by analytic force evaluations \cite{ulibarrena2024hybrid}. In particular, the NN correction
changes only the scalar magnitude of the force; the direction remains the exact radial
Newtonian direction throughout.

To resolve close encounters without globally shrinking the timestep, \texttt{SIMON} uses
adaptive sub-stepping. At each macro step, the minimum pairwise distance $r_{\min}$ is
computed. If $r_{\min}<0.05\,\mathrm{AU}$, the macro step $dt$ is divided into
\[
n_{\mathrm{sub}}
=
\min\left(
16,\,
\max\left(2,\left\lceil \frac{0.05}{r_{\min}}\right\rceil\right)
\right)
\]
sub-steps, each of length $dt/n_{\mathrm{sub}}$. Thus, closer encounters receive finer temporal
resolution, while the maximum of 16 sub-steps limits computational overhead. All positions,
velocities, accelerations, and forces are stored and updated in float64; the neural-network
forward pass uses float32 weights, with the scalar correction converted back to float64 before
being applied.

\subsection{Training Procedure}
\label{sec:training}

Training data is generated from synthetic pairwise configurations sampled across the
separation range relevant to the force correction. The target is the analytic correction
$\log c_{ij}=3\log(\tilde r_{ij}/r_{ij})$, so the network learns how to convert the softened
force magnitude back toward the Newtonian force magnitude while keeping the force direction
analytic. Approximately $60\%$ of samples are drawn from the transition region
$r\in[0.5\epsilon,50\epsilon]$, where softening changes the force appreciably, and $40\%$
from the far field $r\in[50\epsilon,10]$, where the ideal correction is close to one.

\begin{table}[H]
\centering
\caption{Training configuration for the scalar force-correction network.}
\label{tab:training_config}
\begin{tabular}{@{}ll@{}}
\toprule
\textbf{Component} & \textbf{Setting} \\
\midrule
Input features & $[\log \tilde r_{ij},\log m_i,\log m_j]$ \\
Target & $\log c_{ij}=3\log(\tilde r_{ij}/r_{ij})$ \\
Architecture & $3\rightarrow32\rightarrow32\rightarrow32\rightarrow1$ with SiLU activations \\
Training samples & 40,000 \\
Validation samples & 10,000 \\
Optimizer & AdamW, weight decay $1\times10^{-5}$ \\
Learning rate schedule & $10^{-3}$ with cosine annealing to $10^{-5}$ \cite{loshchilov2017sgdr} \\
Loss & Mean squared error on $\log c_{ij}$ \\
Epochs & 5,000 \\
Primary training seed & 42 \\
Training device & NVIDIA GeForce RTX 5050 CUDA 13.0 \\
Best validation MSE & $1.1\times10^{-5}$ \\
\bottomrule
\end{tabular}
\end{table}

Seed robustness is tested separately in Section~\ref{sec:seed_robustness} by retraining the
same architecture with five independent random initialisations.

\subsection{Evaluation Strategy and Scope}
\label{sec:evaluation}

Because three-body trajectories are chaotic, long-horizon pointwise agreement alone is not
a sufficient evaluation standard. \texttt{SIMON} is therefore evaluated using both quantitative
and qualitative diagnostics:

\begin{table}[H]
\centering
\caption{Evaluation diagnostics used to assess \texttt{SIMON}.}
\label{tab:evaluation_strategy}
\begin{tabular}{@{}p{0.30\linewidth}p{0.62\linewidth}@{}}
\toprule
\textbf{Diagnostic} & \textbf{Purpose} \\
\midrule
Trajectory overlays & Visual comparison of \texttt{SIMON} and \texttt{ias15} trajectories, including phase drift. \\
RMS separation & Measures accumulated position deviation from the \texttt{ias15} reference. \\
Divergence-rate proxy $\lambda$ & Summarises early-to-mid horizon separation growth in chaotic rollouts. \\
Bounded/ejected status & Tests whether \texttt{SIMON} preserves the qualitative dynamical regime. \\
Energy drift & Checks whether physical consistency follows the expected leapfrog $dt^2$ scaling. \\
Seed robustness & Tests whether results depend on random NN initialisation. \\
Speed--accuracy frontier & Compares runtime gain against trajectory error across timesteps. \\
\bottomrule
\end{tabular}
\end{table}

The fitted value of $\lambda$ used in this work is a divergence-rate proxy rather than a
formal Lyapunov exponent \citep{sep-chaos}. It reflects both the intrinsic chaotic sensitivity
of the three-body system and the accumulated model/integration mismatch between
\texttt{SIMON} and \texttt{ias15}. A perfect reproduction of the reference trajectory would give
$\lambda=0$ in this comparison, while larger values indicate faster separation from the
reference rollout.

The scope of the method is deliberately limited. \texttt{SIMON} is tested as a controlled
methodological study for unrestricted three-body systems, not as a universal replacement for
\texttt{ias15}. Its purpose is to determine whether a neural force correction can remain useful
when exact radial direction and close-encounter timestep refinement are enforced analytically.
The limitations of this scope are discussed further in Section~\ref{sec:limitations}.

% =============================================================================
\section{Results}
\label{sec:results}
% =============================================================================

In this section we present the experimental results for \texttt{SIMON} across four
themes. Section~\ref{sec:stability} isolates the individual contribution of each
architectural component through a controlled four-run ablation. Section~\ref{sec:generalisation}
evaluates generalisation across seven initial conditions. Section~\ref{sec:physical}
examines physical consistency and reproducibility. Section~\ref{sec:computational_performance}
characterises the speed-accuracy trade-off as a separate computational-performance result.

Throughout this section, $\lambda$ denotes the fitted divergence-rate proxy (slope of
$\log\delta(t)$ in the 10--50\,yr window), and ``bounded'' indicates that all three
bodies remain within 10\,AU of the system barycentre at $T = 100$\,yr. The reference
integrator is \texttt{ias15} \citep{Rein_2014}; all initial conditions use masses
$[1.0,\,0.01,\,0.005]\,M_\odot$ unless otherwise stated. At the operational timestep
$dt = 0.04$\,yr, $\lambda = 0.1655$\,yr$^{-1}$ and final RMS $= 1.54$\,AU for the
default initial condition (IC1).

\subsection{Results I --- Stability Decomposition}
\label{sec:stability}

\subsubsection{The Four-Run Ablation Design}

To isolate the role of each architectural component, we conduct a controlled four-run
ablation in which exactly one component is removed or altered per run while the initial
condition IC1, timestep grid, simulation horizon, and evaluation protocol remain fixed. Runs
A, C, and D use the same trained SIMON weights; Run B uses a separately trained
Vector NN baseline because the architecture itself is changed to learn force direction.
The four runs are summarised in
Table~\ref{tab:ablation}.

\begin{table}[H]
\centering
\caption{Four-run ablation framework for IC1 over $T=100$\,yr. Each run removes or replaces one component of the full \texttt{SIMON} configuration.}
\label{tab:ablation}
\resizebox{\columnwidth}{!}{%
\begin{tabular}{@{}lllll@{}}
\toprule
\textbf{Run} & \textbf{Adaptive stepping} & \textbf{Force direction} & \textbf{NN correction} & \textbf{Outcome} \\
\midrule
A & \textbf{OFF} & Analytic & ON  & Ejection (161.7\,AU) \\
B & ON & \textbf{Learned (Vector NN)} & N/A & Ejection (83.0\,AU) \\
C & ON & Analytic & \textbf{OFF} & Bounded (1.67\,AU) \\
D (\textsc{SIMON}) & ON & Analytic & ON & Bounded (1.54\,AU) \\
\bottomrule
\end{tabular}%
}
\end{table}

\noindent The logic of the design is as follows. Runs A and B establish stability requirements by
showing that removing a component causes ejection. Run C establishes the NN’s role by
showing that its removal leaves the system stable. Run D is the full system, providing
the accuracy ceiling against which Run C is compared. Because Run B is evaluated on the
default initial condition, its trajectory outcome is interpreted together with the
gate-sensitivity analysis rather than as a standalone statistical generalisation.

\subsubsection{Run A: Adaptive Sub-Stepping Is the Primary Stability Mechanism}
\label{sec:run_a}

Without adaptive sub-stepping, \texttt{SIMON} uses a fixed leapfrog step for every
integration step regardless of body separation. At $dt = 0.01$\,yr, a close encounter
occurs whose duration is shorter than the fixed timestep can resolve. The lightest
body receives an incorrectly large gravitational impulse in a single step and is ejected
to $161.7$\,AU at $T = 100$\,yr. With adaptive sub-stepping active, the same close encounter is resolved by splitting the step
into 2 to 16 sub-steps when $r_{\min} < 0.05$ AU, and ejection is eliminated across all tested
timesteps for IC1.

Figure~\ref{fig:run_a_frontier} shows the complete frontier
sweep across five timestep values: $dt = 0.005$, $0.01$, $0.02$,
$0.04$, and $0.08$\,yr. The sweep is run for both conditions. The
worst-case final RMS across the frontier reduces from $161.7$\,AU
before adaptive sub-stepping ($dt = 0.01$\,yr) to $1.54$\,AU after
adaptive sub-stepping ($dt = 0.04$\,yr), a $105\times$ improvement.

\begin{figure}[H]
    \centering
    \includegraphics[width=0.82\linewidth]{adaptive_error_vs_dt.png}
    \caption{Final RMS position error vs.\ timestep for \texttt{SIMON} with and without adaptive sub-stepping on IC1.}
    \label{fig:run_a_frontier}
\end{figure}

At the operational timestep $dt = 0.04$\,yr---the main timestep used for the
operational comparison---adaptive sub-stepping costs $+28\%$ per-step overhead
($19.21 \to 24.57\,\mu$s), reducing speedup from $1.53\times$ to $1.12\times$
relative to \texttt{ias15}, while reducing final RMS from $1.691$ to $1.543$\,AU
($8.8\%$ improvement) and the divergence-rate proxy from $0.1712$ to
$0.1655$\,yr$^{-1}$ ($3.3\%$ improvement). Both configurations remain faster
than \texttt{ias15}. The reported NN invocation rate is less than $1.1\%$ of pair evaluations
involving close encounters with $r<0.15$ AU. Adaptive sub-stepping uses the
stricter condition $r_{\min}<0.05$ AU, so it is triggered only on a subset of
those close-encounter cases, leaving the majority of steps at the base leapfrog
cost.

\subsubsection{Run B: Analytic Force Direction as a Structural Physical Constraint}
\label{sec:run_b}

In Run B, the analytic force direction $\mathbf{u}_{ij} = \mathbf{r}_{ij}/\|\mathbf{r}_{ij}\|$ is replaced by a learned direction from a Vector NN. The Vector NN uses the same hidden-layer structure as \texttt{SIMON} but receives strictly more information, including the geometric direction components $[r_x/r, r_y/r, r_z/r]$, and predicts a unit force-direction vector. Post-training verification confirms sub-degree directional accuracy, with maximum directional error approximately $0.04^\circ$ across the tested distance range. For the tested initial condition, however, the learned-direction architecture produces unbounded trajectories, reaching final RMS error $83.0$\,AU at $T = 100$\,yr, while \texttt{SIMON} remains bounded at $1.54$\,AU. In plain terms, this run tests whether gravity's direction should be learned approximately or enforced exactly from geometry.

To test whether this unbounded outcome was caused by a single anomalous direction prediction or by accumulated small errors, the Vector NN direction gate was tightened and the simulation was re-run.

\begin{table}[H]
\centering
\caption{Gate-sensitivity test for the learned-direction Vector NN baseline. Tightening the direction gate does not prevent the unstable outcome.}
\label{tab:vector_gate_sensitivity}
\begin{tabularx}{\linewidth}{
    >{\raggedright\arraybackslash}p{0.24\linewidth}
    >{\raggedright\arraybackslash}X
    >{\raggedright\arraybackslash}p{0.28\linewidth}
}
\toprule
\textbf{Gate setting} & \textbf{Interpretation} & \textbf{Outcome} \\
\midrule
$\cos(\theta) < 0.866$
&
Standard $30^\circ$ gate; triggers only on large direction deviations.
&
Ejection unchanged; final RMS $83.0$\,AU.
\\
\addlinespace
$\cos(\theta) < 0.9999985$
&
Much tighter than the standard gate; close to the Vector NN's measured sub-degree accuracy.
&
Identical ejection timing; final RMS $83.0$\,AU.
\\
\bottomrule
\end{tabularx}
\end{table}

The gate-sensitivity result shows that the mechanism is not a single catastrophic misprediction, but the accumulation of small non-radial impulses over many steps. When the gate is tightened to $\cos(\theta) < 0.9999985$, corresponding to an angular threshold of approximately $0.1^\circ$, the ejection timing and final RMS remain unchanged. This threshold is far tighter than the original $30^\circ$ gate and close to the Vector NN's measured sub-degree directional accuracy. 

Thus, the failure is not an outlier that a gate can catch. Instead, many individually acceptable direction errors accumulate during close encounters. A simple scale estimate illustrates the issue: at $dt = 0.04$\,yr over $100$\,yr, there are approximately $2{,}500$ integration steps, and $0.04^\circ$ per close-encounter step can accumulate to roughly $100^\circ$ of non-radial impulse. Because Newtonian gravity is a central force directed exactly along the line connecting two bodies, any learned perpendicular component has no physical basis and can compound into non-physical energy injection in a chaotic system.

Figures~\ref{fig:run_b_divergence} and~\ref{fig:run_b_trajectory} illustrate the physical consequence of the gate-sensitivity finding. During the bounded phase, \texttt{SIMON} and the Vector NN have nearly identical fitted divergence rates ($0.1655$ vs.\ $0.1656$\,yr$^{-1}$), but the learned-direction model later becomes unbounded for this tested initial condition, reaching final RMS error $83.0$\,AU. In Figure~\ref{fig:run_b_trajectory}, the long red dotted segments are escape paths of the Vector NN trajectory, not stable orbital loops.

\begin{figure}[H]
    \centering
    \includegraphics[width=0.82\linewidth]{divergence_3model.png}
    \caption{RMS position error vs.\ time for \texttt{SIMON} and the Vector NN baseline. The shaded region marks the 10--50\,yr divergence-rate fit window.}
    \label{fig:run_b_divergence}
\end{figure}

\begin{figure}[H]
    \centering
    \includegraphics[width=0.95\linewidth]{trajectory_3model.png}
    \caption{XY trajectory overlays for \texttt{ias15}, \texttt{SIMON}, and the Vector NN over 100 years. Each body is shown separately with independent axis limits.}
    \label{fig:run_b_trajectory}
\end{figure}

This result supports the conclusion that analytic force direction is not merely a
convenience but a structural stability constraint. For the tested initial condition,
the learned-direction baseline becomes unbounded despite sub-degree angular accuracy,
and the gate-sensitivity analysis shows that the outcome is insensitive to the specific
gate threshold chosen. Analytic direction therefore removes, by construction, a
non-physical perpendicular-force error mode that a learned direction model can still
introduce during close encounters.

\subsubsection{Run C: The Neural Network Correction Is an Accuracy Mechanism}
\label{sec:run_c}

In Run C, the NN correction is disabled entirely: the correction factor $c$ is fixed
at $1.0$ for all pair evaluations, reducing the force to pure softened Newtonian
gravity. The analytic force direction and adaptive sub-stepping remain active. The
result is a bounded simulation at $1.67$\,AU---not an ejection.

Figure~\ref{fig:run_c_divergence} shows that Run C (No-NN baseline, red dashed) and
Run D (full \texttt{SIMON}, green solid) are visually indistinguishable throughout
the 100-year horizon. Both achieve $\lambda = 0.1655$\,yr$^{-1}$ to four decimal
places. The NN correction contributes $8.5\%$ improvement in final RMS accuracy
($1.67 \to 1.54$\,AU) at the $0.45\%$ of pair evaluations involving close encounters
below $0.15$\,AU, where it is invoked.

\begin{figure}[H]
    \centering
    \includegraphics[width=0.92\linewidth]{divergence_no_nn.png}
    \caption{RMS position error vs.\ time for full \texttt{SIMON} and the No-NN baseline. Both runs retain analytic direction and adaptive sub-stepping.}
    \label{fig:run_c_divergence}
\end{figure}

This result clarifies the role of the NN correction. Run A shows that close-encounter
stability is provided by adaptive sub-stepping, since removing that mechanism causes
ejection even when the NN correction remains active. Run C shows that removing the
NN correction while retaining adaptive sub-stepping leaves the system bounded, but
increases final RMS error from $1.54$\,AU to $1.67$\,AU. Therefore, the NN correction
is best interpreted as a targeted accuracy refinement at close-encounter steps, rather
than the primary stability mechanism.

\subsubsection{Synthesis: The Two-Constraint Principle}
\label{sec:synthesis}

The four-run ablation establishes a clear decomposition of \texttt{SIMON}'s performance:

\begin{itemize}
  \item \textbf{Adaptive sub-stepping} eliminates the leapfrog ejection failure at
$dt=0.01$ yr for IC1 (Run A: $161.7$\,AU without sub-stepping versus
$0.847$\,AU with sub-stepping).
  \item \textbf{Analytic force direction} is a structural physical constraint (Run B:
        for the tested initial condition, replacing analytic direction with a learned
        direction produces unbounded trajectories with final RMS $83.0$\,AU).
  \item \textbf{NN correction} is an accuracy mechanism (Run C: removing it leaves
        stability intact but costs $8.5\%$ accuracy).
  \item \textbf{Full \texttt{SIMON}} (Run D) achieves the best accuracy by combining
        all three components.
\end{itemize}

Figure~\ref{fig:simon_traj} shows the XY trajectory overlay and coordinate time series
for the full \texttt{SIMON} system (Run D), confirming bounded dynamics for all three
bodies over 100 years.

\begin{figure}[H]
    \centering
    \includegraphics[width=0.95\linewidth]{traj_overlay_xy_subplots_all_bodies_T100_after_adaptive.png}
    \caption{XY trajectory overlays for all three bodies comparing \texttt{SIMON} with \texttt{ias15} over 100 years.}
    \label{fig:simon_traj}
\end{figure}

\begin{figure}[H]
    \centering
    \includegraphics[width=0.95\linewidth]{traj_overlay_timeseries_all_bodies_T100_after_adaptive.png}
    \caption{Coordinate time series $x(t)$ and $y(t)$ for all three bodies over 100 years. Phase drift accumulates while the trajectories remain bounded.}
    \label{fig:simon_timeseries}
\end{figure}

The primary error mode is phase drift---a gradual timing mismatch that accumulates
while preserving the qualitative orbital geometry and boundedness of all three bodies.
This is consistent with the chaotic sensitivity of the Three-Body Problem
\citep{ott2002chaos} and does not indicate structural failure of the integrator.

% ── Results II ───────────────────────────────────────────────────────────────

\subsection{Results II --- Generalisation Across Initial Conditions}
\label{sec:generalisation}

\subsubsection{Seven Initial Conditions Spanning the Dynamical Spectrum}

To test whether \texttt{SIMON}'s performance extends beyond the single default
configuration (IC1) used in the ablation study, we evaluate across seven distinct
three-body initial conditions designed to span qualitatively different dynamical
regimes. Table~\ref{tab:multi_ic} summarises the configuration and simulation outcomes
for all seven ICs. All runs use masses $[1.0,\,0.01,\,0.005]\,M_\odot$ except IC2
(near-equal mass), $T = 100$\,yr, and $dt = 0.04$\,yr.

\begin{table}[H]
\centering
\caption{Multi-IC evaluation results at $T=100$\,yr and $dt=0.04$\,yr. Summary statistics are computed over the four bounded configurations only.}
\label{tab:multi_ic}
\begin{tabular}{@{}llllll@{}}
\toprule
\textbf{IC} & \textbf{Regime} & $\boldsymbol{\lambda}$ (yr$^{-1}$) & \textbf{RMS (AU)} & \textbf{Bounded} & \textbf{Speedup} \\
\midrule
IC1: Default & Moderate scattering & 0.1655 & 1.54 & Yes & 1.12$\times$ \\
IC2: Near-equal & Strongly interacting & 0.0430 & 116.9 & No & 1.77$\times$ \\
IC3: Tight pair & High encounter rate & 0.1207 & 3.45 & Yes & 1.24$\times$ \\
IC4: Hierarchical & Weakly coupled outer & 0.0990 & 1.60 & Yes & 1.14$\times$ \\
IC5: High ecc.\ (v1) & High eccentricity & 0.0648 & 20.0 & No & 1.22$\times$ \\
IC6: Near-circular & Low chaos & $-0.019$ & 0.008 & Yes & 1.08$\times$ \\
IC7: High ecc.\ (v2) & High eccentricity & 0.0341 & 11.4 & No & 1.11$\times$ \\
\midrule
\multicolumn{2}{l}{Mean (bounded ICs only)} & 0.092 & 1.65 & -- & -- \\
\multicolumn{2}{l}{Std (bounded ICs only)} & 0.068 & 1.22 & -- & -- \\
\multicolumn{2}{l}{Bounded count} & -- & -- & 4/7 & -- \\
\bottomrule
\end{tabular}
\end{table}

\subsubsection{Bounded Configurations: Four ICs}
\label{sec:bounded_ics}

All four stable configurations maintain bounded dynamics for the full 100-year
horizon. The configurations span the complete stable dynamical spectrum of the
three-body problem:

\begin{itemize}
  \item \textbf{IC6 (near-circular)}: $\lambda \approx 0$, final RMS $= 0.008$\,AU.
        The NN is never invoked (NN\_frac $= 0.000$); \texttt{SIMON} computes exact
        Newtonian forces for every step. The residual $0.008$\,AU is the pure
        leapfrog--Gauss-Radau integrator difference.
  \item \textbf{IC4 (hierarchical)}: $\lambda = 0.099$\,yr$^{-1}$, final RMS
        $= 1.60$\,AU. The outer body at $5$\,AU is weakly coupled; the NN is also
        never invoked.
  \item \textbf{IC1 (default)}: $\lambda = 0.1655$\,yr$^{-1}$, final RMS
        $= 1.54$\,AU. The default IC1 baseline.
  \item \textbf{IC3 (tight binary)}: $\lambda = 0.121$\,yr$^{-1}$, final RMS
        $= 3.45$\,AU. The inner body at $0.5$\,AU generates continuous close
        encounters; despite this, \texttt{SIMON} remains bounded throughout.
\end{itemize}

Figure~\ref{fig:multi_ic_results} summarises the divergence rates, final RMS values,
and speedups across all seven ICs.

\begin{figure}[H]
    \centering
    \includegraphics[width=0.95\linewidth]{multi_ic_results_v3.png}
    \caption{Summary across seven initial conditions: divergence rate, final RMS error, and speedup relative to \texttt{ias15}. Hatched bars denote physically unstable configurations.}
    \label{fig:multi_ic_results}
\end{figure}

Figure~\ref{fig:multi_ic_divergence} shows the divergence curves for all seven ICs
simultaneously.

\begin{figure}[H]
    \centering
    \includegraphics[width=0.95\linewidth]{multi_ic_divergence_v3.png}
    \caption{RMS position error vs.\ time for all seven initial conditions. The shaded region marks the 10--50\,yr divergence-rate fit window.}
    \label{fig:multi_ic_divergence}
\end{figure}

\subsubsection{Physically Unstable Configurations}
\label{sec:unstable_ics}

Three configurations produce ejections (IC2, IC5, IC7). In each case, the reference
integrator \texttt{ias15} independently also ejects a body, confirming these are
physical instabilities rather than \texttt{SIMON} failures:

\begin{itemize}
  \item \textbf{IC2 (near-equal mass)}: \texttt{SIMON} ejects to $116.9$\,AU;
        \texttt{ias15} ejects to $58.9$\,AU. The $2\times$ post-ejection discrepancy
        is consistent with chaotic amplification of integrator differences after the
        system becomes dynamically unbound.
  \item \textbf{IC5 (high eccentricity, $v = 2.2$, $r = 0.3$\,AU)}: Both
        integrators eject; \texttt{ias15} ejects Body~2 at $t = 34.8$\,yr.
  \item \textbf{IC7 (high eccentricity redesign, $v = 1.6$, $r = 0.6$\,AU)}:
        Both integrators eject; \texttt{ias15} ejects Body~2 at $t = 88.4$\,yr.
\end{itemize}

\texttt{SIMON} does not suppress these physical instabilities. This is expected and physically consistent: the model agrees with \texttt{ias15} on the occurrence of dynamical instability rather than artificially stabilising configurations that
are physically unbound. The post-ejection distances are not expected to match exactly
because chaotic divergence amplifies integrator differences after unbinding. High-eccentricity
configurations with these mass ratios proved intrinsically unstable in both cases
tested: the eccentric inner body transfers orbital energy to the outer body through
repeated periapsis passages, eventually unbinding it regardless of integrator.

% ── Results III ──────────────────────────────────────────────────────────────

\subsection{Results III --- Physical Consistency and Reproducibility}
\label{sec:physical}

\subsubsection{Energy Conservation: Leapfrog Attribution}
\label{sec:energy}

We track fractional energy drift $\Delta E(t)/E_0 = (E(t) - E_0)/|E_0|$ throughout
each simulation, where $E(t) = \mathrm{KE}(t) + \mathrm{PE}(t)$ is the total
mechanical energy. \texttt{ias15} conserves energy to float64 machine precision
($\approx 10^{-14}\%$). \texttt{SIMON} at the operational timestep ($dt = 0.04$\,yr)
shows:

\begin{itemize}
  \item IC1: $1.39\%$ maximum fractional drift (NN\_frac $= 0.0045$)
  \item IC3: $0.67\%$ maximum fractional drift (NN\_frac $= 0.000$)
  \item IC4: $13.77\%$ maximum fractional drift (NN\_frac $= 0.000$)
\end{itemize}

IC3 and IC4 both have NN\_frac $= 0.000$: the NN correction was never applied during
their simulations. Yet both show comparable energy drift to IC1. The NN correction
cannot cause energy drift it never contributed to; the drift must arise from the
leapfrog integrator. This is the attribution proof.

IC4's higher fractional drift ($13.77\%$) reflects its small binding energy
$|E_0| = 0.006$\,(AU$^2$/yr$^2$)---the smallest of the three ICs tested. Because
fractional drift is normalised by $|E_0|$, weakly bound hierarchical configurations
can show larger relative energy drift at the same timestep. This should not be
interpreted as NN-driven energy injection: the NN was never invoked for IC4
(NN\_frac $= 0.000$). Importantly, the drift is oscillatory throughout---not
monotonically increasing---consistent with symplectic shadow-Hamiltonian
behaviour, and all bodies remain bounded (IC4 final RMS $= 1.60$\,AU).

Figure~\ref{fig:energy_drift} shows the energy drift time series for all three ICs.

\begin{figure}[H]
    \centering
    \includegraphics[width=0.95\linewidth]{energy_drift_all_ics.png}
    \caption{Fractional energy drift $|\Delta E/E_0|$ over 100 years for \texttt{ias15} and \texttt{SIMON} at $dt=0.04$\,yr across three initial conditions.}
    \label{fig:energy_drift}
\end{figure}

\subsubsection{Timestep Scaling: Quantitative Confirmation of $dt^2$ Prediction}
\label{sec:dt_scaling}

Second-order leapfrog theory predicts that energy drift should scale as $dt^2$:
halving the timestep should reduce drift by a factor of four. We test this by
running \texttt{SIMON} on IC1, IC3, and IC4 at two timestep values ($dt = 0.01$\,yr
and $dt = 0.04$\,yr). The fitted log-log slopes are:

\begin{itemize}
  \item IC1: slope $= 1.32$ (NN\_frac $= 0.0045$)
  \item IC3: slope $= 2.43$ (NN\_frac $= 0.000$)
  \item IC4: slope $= 1.86$ (NN\_frac $= 0.000$)
  \item Average: $1.87$ (theoretical prediction: $2.0$)
\end{itemize}

The average slope of $1.87$ is consistent with $dt^2$ scaling, given that only two
data points are available per IC. For IC4 specifically, the drift reduces from
$13.77\%$ at $dt = 0.04$\,yr to $1.05\%$ at $dt = 0.01$\,yr---a factor of
$13.2\times$, close to the $dt^2$ prediction of $16\times$. IC3 and IC4 follow this scaling despite NN\_frac $= 0.000$ at both timesteps,
confirming that their energy drift is controlled by the leapfrog timestep rather
than by the neural correction.

Figure~\ref{fig:dt_scaling} shows the log-log scaling plot.

\begin{figure}[H]
    \centering
    \includegraphics[width=0.82\linewidth]{01_drift_scaling_loglog.png}
    \caption{Maximum fractional energy drift vs.\ timestep for IC1, IC3, and IC4. The grey reference line shows the expected $dt^2$ leapfrog scaling.}
    \label{fig:dt_scaling}
\end{figure}

\subsubsection{Training Seed Robustness}
\label{sec:seed_robustness}

To verify that results from the seed-42 model used throughout this paper are
not specific to a particular random initialisation, training was repeated five times
with seeds $\{42, 123, 456, 789, 1000\}$. Table~\ref{tab:seeds} summarises the
training and simulation outcomes.

\begin{table}[H]
\centering
\caption{Training seed robustness for the full \texttt{SIMON} run on IC1 at $T=100$\,yr and $dt=0.04$\,yr.}
\label{tab:seeds}
\begin{tabular}{@{}llll@{}}
\toprule
\textbf{Seed} & \textbf{Val.\ MSE} & $\boldsymbol{\lambda}$ (yr$^{-1}$) & \textbf{Bounded} \\
\midrule
42 (primary model) & $1.11\times10^{-5}$ & 0.1655 & Yes \\
123                & $1.72\times10^{-5}$ & 0.1655 & Yes \\
456 (worst)        & $1.75\times10^{-5}$ & 0.1655 & Yes \\
789                & $1.68\times10^{-5}$ & 0.1655 & Yes \\
1000               & $1.43\times10^{-5}$ & 0.1655 & Yes \\
\midrule
Mean               & $1.54\times10^{-5}$ & 0.1655 & 5/5 \\
CV                 & $15.8\%$            & $0\%$  & --- \\
\bottomrule
\end{tabular}
\end{table}

The divergence rate $\lambda = 0.1655$\,yr$^{-1}$ is identical to four decimal places across all five seeds. All five independently trained models produce bounded dynamics with identical fitted divergence rate, confirming that downstream stability is insensitive to random initialisation.

The $15.8\%$ coefficient of variation in validation MSE does not imply dynamical
sensitivity to seed choice: the variation is concentrated at sub-softening distances
($r < 5\times10^{-4}$\,AU) where the safety gate bypasses NN predictions entirely.
In the operationally relevant range ($r \geq 10^{-3}$\,AU), all seeds produce
correction factor errors below $0.3\%$---functionally identical---and the simulation
metrics are invariant.

Figure~\ref{fig:seeds} shows the training convergence and per-seed validation MSE.

\begin{figure}[H]
    \centering
    \includegraphics[width=0.95\linewidth]{seed_robustness_results.png}
    \caption{Training seed robustness across five random initialisations. Left: validation MSE curves; right: best validation MSE per seed.}
    \label{fig:seeds}
\end{figure}

% ── Results IV ───────────────────────────────────────────────────────────────

\subsection{Results IV --- Computational Performance}
\label{sec:computational_performance}

\subsubsection{Speed-Accuracy Frontier}
\label{sec:perf}

In this measurement, we observe the final position deviation of the bodies predicted
by \texttt{SIMON}, in comparison with \texttt{ias15}, over a 100 year time horizon.
The x-axis is the throughput, which is how often the current state of the system is
sampled, \texttt{Rebound}, and the y-axis is RMS position error, on a logarithmic
scale. The total number of samples is fixed across all simulations. The timestep $dt$
is the timestep size the hybrid model uses internally. Points toward the lower-right
of the graph are better (higher throughput, lower error). Throughput is computed as
$\frac{n_{\mathrm{samples}}}{t_{\mathrm{total}}}$.

\begin{figure}[H]
    \centering
    \includegraphics[width=0.82\linewidth]{speed_accuracy_frontier_T100_after_adaptive.png}
    \caption{Speed--accuracy frontier for \texttt{SIMON} across timestep values. Lower-right points indicate higher throughput and lower final RMS error.}
    \label{fig:frontier}
\end{figure}

The \texttt{ias15} baseline completes in $0.080$ seconds. \texttt{SIMON} outperforms
\texttt{ias15} at $dt = 0.04$ ($1.12\times$ speedup) and $dt = 0.08$ ($2.18\times$
speedup). The lowest final RMS error is $0.847$ at $dt = 0.01$. This is the optimal
point for accuracy, but it is approximately $3.4\times$ slower than \texttt{ias15}.
The errors at $dt = 0.005$, $0.02$, and $0.08$ are all approximately $1.02$,
reflecting the chaotic saturation scale. After 100 years, any two different integrations
of the same system will diverge to roughly one orbital radius, regardless of timestep.
$dt = 0.04$ shows a slightly higher error of $1.54$, which is a chaotic fluctuation
(a different phase of the orbit was reached at $T = 100$).

Table~\ref{tab:frontier} provides the full quantitative results for the frontier sweep.

\begin{table}[H]
\centering
\caption{Speed--accuracy frontier results for \texttt{SIMON} over $T=100$\,yr. The bold row marks the operational timestep.}
\label{tab:frontier}
\resizebox{\columnwidth}{!}{%
\begin{tabular}{@{}lllllll@{}}
\toprule
$dt$ (yr) & Final RMS (AU) & Throughput (samp/s) & Time/step ($\mu$s) & Total time (s) & Steps & Speedup \\
\midrule
$0.005$  & $1.022$ & $9{,}285$   & $26.9$ & $0.538$ & $20{,}000$ & $0.15\times$ \\
$0.010$  & $0.847$ & $18{,}307$  & $27.3$ & $0.273$ & $10{,}000$ & $0.29\times$ \\
$0.020$  & $1.022$ & $35{,}672$  & $28.0$ & $0.140$ & $5{,}000$  & $0.57\times$ \\
$\mathbf{0.040}$ & $\mathbf{1.543}$ & $\mathbf{69{,}880}$ & $\mathbf{28.6}$ & $\mathbf{0.072}$ & $\mathbf{2{,}500}$ & $\mathbf{1.12\times}$ \\
$0.080$  & $1.022$ & $135{,}629$ & $29.5$ & $0.037$ & $1{,}250$  & $2.18\times$ \\
\bottomrule
\end{tabular}%
}
\end{table}

The per-step time ranges from $26.9\,\mu$s at $dt = 0.005$ to $29.5\,\mu$s at
$dt = 0.08$: a range of only $2.6\,\mu$s across a $16\times$ change in timestep.
This nearly flat profile indicates that the cost of an individual \texttt{SIMON}
step is approximately constant and is dominated by force evaluation, pairwise
distance checks, and occasional adaptive sub-stepping. Therefore, the observed
runtime improvement at larger $dt$ comes primarily from taking fewer macro steps,
not from making each individual step substantially cheaper. For a fixed
$T = 100$\,yr horizon, the number of macro steps scales approximately as $T/dt$,
so total runtime scales roughly as $(T/dt)$ times the near-constant per-step cost.

% =============================================================================

\section{Discussion and Conclusion}
\label{sec:discussion}
% =============================================================================

\subsection{The Two-Constraint Principle}

The central finding of this work is that long-horizon orbital stability in a hybrid
neural three-body integrator is not produced by the neural network alone. It requires
two complementary physical constraints: adaptive sub-stepping and analytic force
direction. Adaptive sub-stepping operates at the integration level by resolving close
encounters that a fixed macro-timestep can under-resolve. Analytic force direction
operates at the force-computation level by enforcing the exact Newtonian radial
symmetry at every step.

The ablation study shows that these constraints have distinct roles. Removing adaptive
sub-stepping causes an under-resolved close encounter to produce ejection in the timestep
frontier. Replacing analytic force direction with a learned direction also leads to
unbounded motion for the tested initial condition. By contrast, removing the neural
correction while retaining both physical constraints preserves bounded dynamics, with
only a modest increase in final RMS error. Thus, the neural network correction is best
understood as an accuracy refinement, while adaptive sub-stepping and analytic direction
are the main stability mechanisms.

The broader implication is that neural models for chaotic physical systems should not
freely learn quantities whose symmetries are already known exactly. In \texttt{SIMON},
Newtonian radial direction is enforced analytically, and the network is restricted to a
low-dimensional scalar correction during close encounters. This suggests a general design
principle for hybrid scientific machine learning: preserve exact physical structure where
available, and use learning only for targeted corrections where analytic or numerical
approximations are most fragile.

\subsection{Positioning \texttt{SIMON} Relative to \texttt{ias15}}

\texttt{SIMON} should not be interpreted as a replacement for \texttt{ias15}. The reference
integrator remains substantially more accurate in energy conservation because it uses
high-order adaptive Gau\ss--Radau integration. \texttt{SIMON}, by contrast, uses a
second-order leapfrog backbone with targeted sub-stepping, so its energy drift is larger
and follows the expected timestep-dependent behaviour of the underlying integrator.

The relevant comparison is therefore not whether \texttt{SIMON} reproduces every
\texttt{ias15} coordinate exactly over 100 years. In a chaotic three-body system, small
model or numerical differences naturally accumulate as phase drift. The more meaningful
question is whether the hybrid method preserves the qualitative dynamical regime:
bounded configurations should remain bounded, physically unstable configurations should
be identified as unstable, and energy drift should remain physically interpretable.

Under this interpretation, \texttt{SIMON} identifies a stability-preserving operating point.
At the operational timestep $dt=0.04$\,yr, it preserves bounded dynamics across all tested
stable configurations while providing a modest computational-cost reduction relative to
\texttt{ias15}. Coarser timesteps can give larger speedups, but they are best interpreted
as part of the speed--accuracy frontier rather than as the main operating regime.

To answer the research question:
\begin{quote}
\emph{How closely can physically constrained neural network force models approximate the
long-term stability of traditional numerical integrators in three-body simulations, while
reducing computational cost without sacrificing qualitative dynamical stability?}
\end{quote}

\noindent In the tested three-body regimes, physically constrained neural force correction
can approximate the qualitative long-horizon behaviour of traditional numerical
integration. All configurations that remain bounded under \texttt{ias15} also remain
bounded under \texttt{SIMON}, and the configurations that eject under \texttt{ias15} are
also identified as physically unstable by \texttt{SIMON}. The main conclusion is therefore
not that \texttt{SIMON} replaces \texttt{ias15}, but that neural force correction can be made
dynamically usable in chaotic gravitational simulation when exact physical structure is
enforced analytically.

\subsection{Limitations and Future Work}
\label{sec:limitations}

This work should be interpreted as a controlled methodological study rather than as a
universal neural replacement for high-order reference integrators. Several limitations
define the scope of the present results and suggest natural directions for future work.

First, the ablation study in Section~\ref{sec:stability} is centred on IC1. Although the
multi-IC evaluation in Section~\ref{sec:generalisation} provides evidence that the full
\texttt{SIMON} configuration generalises across several regimes, the causal decomposition
of adaptive sub-stepping, analytic direction, and neural correction should be repeated over
a broader set of mass ratios and initial geometries.

Second, the fitted divergence-rate $\lambda$ is a proxy rather than a formal Lyapunov
exponent. It compares \texttt{SIMON} against \texttt{ias15}, so it reflects both intrinsic
chaotic divergence and accumulated model mismatch. Future work could compute formal
Lyapunov exponents by evolving nearby initial conditions under the same dynamics and
comparing them with the model-to-reference divergence used here \citep{sep-chaos}.

Third, the operating timestep and close-encounter thresholds are tuned for the tested
three-body systems. Larger systems, different mass ratios, or more extreme eccentricities
may require different thresholds, broader training distributions, or additional safeguards.
The pairwise structure of the correction model should extend naturally to $N>3$ systems,
but this has not yet been demonstrated.

Fourth, the present implementation is not optimised for maximum computational speed.
The measured speedup is modest because runtime includes force evaluation, gating,
sub-stepping, and framework overhead. Future implementations could improve throughput
through vectorised pair evaluation, batching, GPU inference, or lower-level compiled
kernels.

Taken together, these limitations do not weaken the central result; they define its proper
scope. \texttt{SIMON} demonstrates that learned force correction can be useful in chaotic
gravitational simulation when the physically dangerous parts of the problem---force
direction and close-encounter resolution---are constrained analytically rather than left
to an unconstrained neural approximation.

% =============================================================================
% REFERENCES
% =============================================================================
\newpage
\bibliographystyle{plainnat}
\bibliography{V3_bib}


% =============================================================================
\end{document}
% =============================================================================
